\section{Relation to the weak-groupoid expansion method}
\label{sec:context}

The elementary-expansion theorem of \cite{OpenAI2026} concerns
models of a fixed Grothendieck coherator in the Ara--Henry convention.
Its Theorem~1.1 asserts that, for a cellular model \(X\), freely adjoining
a parallel \(m\)-cell and an \((m+1)\)-cell joining it to a prescribed
cell preserves components and based homotopy groups.  Cellular inputs
include transfinite cellular presentations; its Corollary~5.1 extends
the assertion to cofibrant retracts and gives a left semi-model
structure.  These hypotheses differ from those of
Theorem~\ref{thm:couniversal}, whose attaching target is an arbitrary
strict \(n\)-category.

The construction in \cite[Sections~3--4]{OpenAI2026}
first reflects a model to one with unique fillers above a fixed
dimension \(q\), preserving all cells through \(q\).  It then descends
in the expansion dimension.  Within one step, higher expansions
already covered by the induction control upper cellular extensions.
Finite tree diagrams are compared with strings using compositions
and inverses.  Duplicating the top generators and adjoining higher
joining cells promotes the resulting generator homotopies to a
statement about the entire free model.  Partial cylinders, stationary
below the working dimension, are then completed along their upper
boundaries.  Finally, a
homotopy-group test uses only finitely many dimensions and therefore
passes out of the cutoff.  Filtered-colimit closure also handles
transfinite cellular presentations.

There is a formal strict analogue of the cutoff: retain the cells of
a strict \(\omega\)-category through dimension \(q\), and above \(q\)
put one cell over each parallel boundary.  Higher identities and
compositions are forced by their boundaries, and all higher laws
follow from uniqueness.  This construction preserves the
\(q\)-cells, whereas the reflection \(\tau_q\) used in this paper
can identify them.  The formal analogy supplies no comparison with
Street weak equivalences.  For example, the inclusion of the object
\(0\) in the ordinary category \(0\longrightarrow1\longleftarrow2\)
is a Street weak equivalence, since the category has a terminal
object.  There is nevertheless no single arrow joining \(0\) and
\(2\) in either direction.  Thus the weak-groupoid test by exact
boundaries and single joining cells cannot serve as a test for
Street weak equivalences.  Inverse traversal in the tree comparison
also requires an independent replacement in the directed setting.

In the present directed setting, the corresponding role is played by
the positive relative prism of
Lemma~\ref{lem:relative-prism}.  It is stationary on the old sieve,
and its chain formula gives compatible strict maps on whole
right Gray cylinders.  Lemma~\ref{lem:internal-truncation} passes
these cylinders to finite \(n\). When their endpoints fix the old
sieve, Lemma~\ref{lem:pushout-cylinder} glues them at arbitrary attaching
targets. A greatest new element gives two local comparisons of this form;
finite sieve covers glue their weak-equivalence conclusions.
Finite subposets stable under the retraction and filtered colimits
then remove the finiteness restriction.  The resulting attachment
comparison applies to subdivided simplicial boundaries and horns
and yields both unrestricted weak factorization systems.
No result of \cite{OpenAI2026} is a premise of this proof.
